% ===================================================================
\section{Vibecoding: AI-Assisted Lean Development}
\label{ch:methodology}
% ===================================================================

\subsection{What is Vibecoding?}

``Vibecoding'' denotes iterative AI-assisted code or proof generation in which
the human operator specifies intent and reviews output while a generative AI model
proposes syntactic content. In conventional software development, vibecoding lacks
a correctness oracle.

Lean changes this dynamic fundamentally. Lean's type-checker is a deterministic
correctness oracle: it either accepts or rejects every proof. The result is that
vibecoded Lean proof attempts are immediately and unambiguously verified or
falsified by the kernel, independent of the operator's understanding.

This produces a principled workflow:

\begin{center}
\begin{tabular}{l}
\textit{Human}: specifies intent and theorem statement \\
$\downarrow$ \\
\textit{AI}: generates candidate proof or tactic block \\
$\downarrow$ \\
\textit{Lean kernel}: accepts or rejects \\
$\downarrow$ \\
\textit{Human}: routes to next iteration or accepts result
\end{tabular}
\end{center}

When Lean accepts, the result is formally guaranteed, regardless of whether the
human operator fully understands each internal tactic step. This is the key
property that makes vibecoding viable for formal mathematics at scale.

% -------------------------------------------------------------------
\subsection{The FRFP Vibecoding Workflow}
% -------------------------------------------------------------------

The formalization of FRFP followed a six-phase workflow:

\begin{description}
  \item[\textbf{Phase 1: Foundation setup}]
        Install Lean~4 and Mathlib. Define core inductive types (\texttt{Object},
        \texttt{Primitive}) and basic kernel theorems. Establish import structure
        and naming conventions.
        \emph{Human role}: architectural decisions, naming standards, dependency
        graph.

  \item[\textbf{Phase 2: Axiom seeding}]
        AI generates candidate axioms from informal paper statements. Human
        screens for consistency and interpretation accuracy. Lean kernel validates
        typing of all axiom signatures.

  \item[\textbf{Phase 3: Theorem attempts}]
        Human specifies target theorem. AI generates proof attempt (tactic blocks,
        \texttt{by} expressions, \texttt{simp} calls). Lean evaluates; accepted
        proofs are committed. Rejected proofs enter repair loop.

  \item[\textbf{Phase 4: Axiom reduction}]
        Systematic audit of all axioms: categorize as derivable theorem, standard
        reference, or true premise. 41 axioms converted to theorems; all remaining
        axioms cite published references.

  \item[\textbf{Phase 5: Documentation integration}]
        AI cross-references Lean code with paper sections and generates summary
        documents. Human reviews for accuracy.

  \item[\textbf{Phase 6: Final validation}]
        Full clean build (\texttt{lake build}) with zero errors. Axiom audit
        confirms 0 uncited axioms. \texttt{grep sorry} confirms 0 live sorry.
\end{description}

% -------------------------------------------------------------------
\subsection{Quantitative Results}
% -------------------------------------------------------------------

\subsubsection*{Productivity}

\begin{center}
\begin{tabular}{ll}
\toprule
Metric & Value \\
\midrule
Total theorems proven & 269 \\
Average time per theorem (with AI) & 15--30 minutes \\
Average time per complex theorem & 2--4 hours \\
Estimated time without AI & 5--10$\times$ longer \\
Developer count & 1 \\
Total timeline & 6 months \\
Total Lean source lines & $\sim$15,000 \\
\bottomrule
\end{tabular}
\end{center}

\subsubsection*{AI Attempt Success Rates}

\begin{center}
\begin{tabular}{ll}
\toprule
Outcome & Rate \\
\midrule
First attempt accepted by Lean & $\sim$30\% \\
Accepted after 1--2 iterations & $\sim$60\% \\
Required significant human intervention & $\sim$10\% \\
\bottomrule
\end{tabular}
\end{center}

\subsubsection*{Error Categories}

\begin{center}
\begin{tabular}{lll}
\toprule
Error type & Catch rate & Catcher \\
\midrule
Type errors, wrong signatures & $\sim$80\% & Lean kernel immediately \\
Logic errors in proof structure & $\sim$15\% & AI review in next iteration \\
Semantic errors (wrong theorem formalized) & $\sim$5\% & Human review \\
\bottomrule
\end{tabular}
\end{center}

% -------------------------------------------------------------------
\subsection{Failure Mode Taxonomy}
% -------------------------------------------------------------------

Six failure modes were observed repeatedly throughout the project.

\begin{description}
  \item[\textbf{FM-1: Hallucinated Mathlib lemma names.}]
        AI generates tactic blocks that call non-existent Mathlib lemmas. Lean
        rejects immediately with ``unknown identifier.''
        \emph{Mitigation}: provide a local lemma search result before requesting
        the proof.

  \item[\textbf{FM-2: Tactic drift across iterations.}]
        After several repair iterations, the AI may introduce new errors while
        fixing old ones.
        \emph{Mitigation}: reset to the last working state and restart repair.

  \item[\textbf{FM-3: Wrong theorem formalization.}]
        AI produces a formally correct Lean statement that does not match the
        intended informal claim. Lean accepts it, but the statement is weaker.
        \emph{Mitigation}: always write the informal statement as a comment above
        the theorem and compare before committing.

  \item[\textbf{FM-4: Axiom proliferation.}]
        AI may generate proofs that succeed by introducing a new \texttt{axiom}
        rather than deriving from existing structure.
        \emph{Mitigation}: prohibit the \texttt{axiom} keyword in AI-generated
        proof attempts unless explicitly authorized.

  \item[\textbf{FM-5: Import path errors.}]
        AI may generate incorrect \texttt{import} statements.
        \emph{Mitigation}: maintain a canonical import list and include it in
        each AI prompt.

  \item[\textbf{FM-6: Type universe mismatches.}]
        AI places structures in the wrong Lean universe (\texttt{Prop} vs
        \texttt{Type}) requiring cascade edits.
        \emph{Mitigation}: fix universe assignments early in each module.
\end{description}

% -------------------------------------------------------------------
\subsection{What the Kernel Cannot Catch}
% -------------------------------------------------------------------

Two categories remain outside the kernel's reach:

\begin{enumerate}
  \item \textbf{Wrong axiom semantics}: An axiom can be formally well-typed but
        semantically wrong. Human review of axiom interpretations is essential.
  \item \textbf{Missing theorems}: Lean will not tell you that you forgot to
        prove something important.
\end{enumerate}

% -------------------------------------------------------------------
\subsection{The Principal Triad}
% -------------------------------------------------------------------

\begin{center}
\begin{tabular}{lll}
\toprule
Role & Agent & Responsibility \\
\midrule
Semantic controller & Human & Decides what to prove and validates intent \\
Proof generator & AI (LLM) & Proposes tactic blocks and proof terms \\
Correctness verifier & Lean kernel & Accepts or rejects every proof step \\
\bottomrule
\end{tabular}
\end{center}

No single agent in the triad is sufficient on its own.

% -------------------------------------------------------------------
\subsection{Lessons for Future Projects}
% -------------------------------------------------------------------

\begin{enumerate}
  \item \textbf{Establish naming conventions before writing any proofs.} Changing
        conventions later causes cascading edits.
  \item \textbf{Audit axioms in batches, not continuously.} Continuous axiom
        cleanup interrupts proof flow.
  \item \textbf{Maintain a proof completion checklist.} A running list of ``not
        yet proven'' theorems prevents false impressions of completeness.
  \item \textbf{Keep AI context small and focused.} Long prompts cause AI to
        hallucinate plausible-but-wrong lemma names.
  \item \textbf{Accept sorry-first, prove-later.} Using \texttt{sorry}
        placeholders to scaffold proof structure before completion was
        consistently more efficient.
\end{enumerate}
